\section{Lecture 1}

\subsection*{Review of AR(1) process.}
Let \(\epsilon_t \sim\) iid. Let \(a \in \mathbb{R}\) be fixed.
\[X_t = a X_{t-1} + \epsilon_t\]
Unraveling this, we are led to a solution of the recursion
\[X_t = \sum_{i=0}^\infty a^i \epsilon_{t-i}\]
In what sense does \(\sum_{i=0}^\infty a^i \epsilon_{t-i}\) exist? Suppose \(\Var[\epsilon_t] < \infty\). Then the series above converges in the mean-square sense. For any \(\epsilon > 0\),
\[\left\|\sum_{i=N}^M a^i \epsilon_{t-i} \right\|_2 = \sigma^2 \sum_{i=N}^M a^i < \epsilon\]
if \(N, M\) are sufficiently large. This is the Cauchy criterion for the space of square-integrable random variables, and hence a limit exists by completeness. Next, does \(\sum_{i=0}^\infty a^i \epsilon_{t-i}\) converge almost surely?

\begin{center}\rule{0.5\linewidth}{0.5pt}\end{center}

\subsection*{Kolmogorov Two Series Theorem}
Let \(X_n\) be independent random variables. Then \(\sum_{n=1}^\infty X_n\) converges almost surely if \(\sum_{n=1}^\infty \E[X_n]\) and \(\sum_{n=1}^\infty \Var[X_n]\) converge almost surely.

\begin{center}\rule{0.5\linewidth}{0.5pt}\end{center}

\subsection*{Kolmogorov Three Series Theorem}
Let \(X_n\) be independent random variables. Then \(\sum_{n=1}^\infty X_n\) converges almost surely if there exists \(A > 0\) and only if for any \(A > 0\):

\begin{itemize}
\tightlist
\item
  \(\sum_{n=1}^\infty \Pr(|X_n| \geq A)\) converges
\item
  Let \(Y_n = X_n \one[|X_n| \leq A]\). Then

  \begin{itemize}
  \tightlist
  \item
    \(\sum_{n=1}^\infty \E[Y_n] < \infty\)
  \item
    \(\sum_{n=1}^\infty \Var[Y_n] < \infty\)
  \end{itemize}
\end{itemize}

\begin{center}\rule{0.5\linewidth}{0.5pt}\end{center}

Suppose \(\Var[\epsilon_t] < \infty\). We can use the two-series theorem to prove sufficient conditions on \(a\) for which \(\sum_{i=0}^\infty a^i \epsilon_{t-i}\) converges.
\[\sum_{i=0}^\infty \E[a^i \epsilon_{t-i}] = \sum_{i=0}^\infty 0 = 0\]

\[\sum_{i=0}^\infty \Var[a^i \epsilon_{t-i}] = \sum_{i=0}^\infty a^{2i}\sigma^2 < \infty\]

under the assumption that \(a^{2i} < 1\) or equivalently \(|a| < 1\).

Now suppose only that \(E \log |\epsilon_t| < \infty\). This is a very weak moment condition (finite geometric mean). We can use the Kolmogorov three-series theorem to prove that \(\sum_{i=0}^\infty a^i \epsilon_{t-i}\) converges almost surely.

We have written the solution to the AR(1) recursion as a \emph{causal} process:
\[X_t = f(\epsilon_t, \epsilon_{t-1}, \dots),\]
where \(f(x_0, x_1, \dots) = \sum_{i=0}^\infty a^i x_i\).

\textbf{Lemma}

Let \(f(x_0, x_1, \dots)\) be a measurable function and \(\epsilon_t\) an iid sequence. Then
\[X_t = f(\epsilon_t, \epsilon_{t-1}, \dots)\]
defines a stationary process. We call this a \emph{Bernoulli shift process}.

\textbf{Proof.}

The marginal distributions are equal; since \((\epsilon_{t_1}, \epsilon_{t_1 - 1}, \dots)\) is equal in distribution to \((\epsilon_{t_1 + \ell}, \epsilon_{t_1 + \ell - 1}, \dots)\), it follows that \(X_{t_1}\) is equal in distribution to \(X_{t_1 + \ell}\). Next, the joint distribution \((X_{t_1}, X_{t_2})\) is equal to that of \((X_{t_1 + \ell}, X_{t_2 + \ell})\). Assume WLOG that \(t_1 > t_2\). Then
\[X_{t_1} = f(\epsilon_{t_1}, \dots, \epsilon_{t_2 + 1} \epsilon_{t_2}, \epsilon_{t_2 - 1}, \dots)\]

\[X_{t_2} = f(\epsilon_{t_2}, \epsilon_{t_2 - 1}, \dots)\]

So the joint distribution is uniquely characterized by the number of innovations not jointly shared, i.e. \(t_1 - t_2\). Thus the joint distribution is invariant to lag shifts.

\subsection*{Non-Causal AR(1) Process}
Suppose \(|a| > 1\). There is a stationary solution to the AR(1) recursion which is not causal. The AR(1) recursion is equivalent to
\[X_{t-1} = X_t /a - \epsilon_t / a\]
This is an AR(1) recursion forward in time, instead of backwards. The convergence results from above imply that for \(|a| > 1\), the random series
\[X_t = \sum_{i=1}^\infty \frac{(-1)^i}{a^i} \epsilon_{t+i}\]
converges almost surely (or in various moments) to a solution of the AR(1) recursion. Yet this it not a causal solution because it depends on future innovations. This solution is practically useless.

\subsection*{Unit-Root AR(1) Process}
Suppose \(a = 1\). There is no stationary solution to the AR(1) recursion in this case. Attempting to unravel the AR(1) recursion, we get
\[X_t = \epsilon_t + \epsilon_{t-1} + \cdots = \sum_{i=0}^\infty \epsilon_{t-i}\]
Unless \(\epsilon_t = 0\) almost surely, this does not converge to a well-defined random variable. Let \(A > 0\) be arbitrary. The first condition of Kolmogorov's three series theorem necessitates that \(\sum_{i=0}^\infty \Pr(|\epsilon_{t-i}| > A) < \infty\). Since \(\epsilon_{t-i}\) are identically distributed, this necessitates \(\Pr(|\epsilon_t| > A) = 0\). Since \(A > 0\) was arbitrary, this implies \(\epsilon_t = 0\) almost surely.

For any constant \(c\), we can set \(X_t = c\) for \(t \leq 0\) and let \(X_t\) process via the aforementioned AR(1) recursion for \(t \geq 0\). This process is called a \emph{random walk}. It does not forgo its initialization, since \(\E[X_t] = c\) for any \(t\) and \(\Var[X_t] = t \sigma^2\) for \(t \geq 0\), assuming necessary moments.

Suppose \(a = -1\). Then TODO

\subsection*{Nonlinear AR(1) Process}
Let \(f\) be a measurable function. Define
\[X_t = f(X_{t-1}, \epsilon_t)\]
This is the generic form for Markov chain with state space \(\R\). We can unravel the recursion as
\[X_t = f(f(X_{t-2}, \epsilon_{t-1}), \epsilon_t)\]

\[X_t = f(f(f(X_{t-3}, \epsilon_{t-2}), \epsilon_{t-1}), \epsilon_t)\]

We want to impose conditions under which we can define the limit \(g\) of the iterated function system so that
\[X_t = g(\epsilon_t, \epsilon_{t-1}, \epsilon_{t-2}, \dots)\]
The above recursion gives an implicit definition of \(g\) via
\[g(\epsilon_t, \epsilon_{t-1}, \dots) = f(g(\epsilon_{t-1}, \epsilon_{t-2}, \dots), \epsilon_t)\]

\subsection*{Stochastic Gradient Descent}
We want to produce an estimate \(\wh\theta\) of a parameter \(\theta\) of a distribution or system. For instance, if \(f_\theta\) is the marginal likelihood of our data, and we want to perform online MLE. Let \(\wh\theta_0\) be arbitrary and define
\[\wh\theta_t = \wh\theta_{t-1} - \gamma_t \frac{d}{d\theta} \log f_\theta(\epsilon_t) \Big\vert_{\theta = \wh\theta_{t-1}}\]
where \(\gamma_t > 0\) is the learning rate sequence. The upshot of this approach is that we only need one gradient computation per step. This defines a discrete-time Markov process. This Markov chain is homogeneous if \(\gamma_t \equiv \gamma\) is constant.

\subsection*{AR(p) Process}
Consider the recursion
\[X_t = a_1 X_{t-1} + \cdots + a_p X_{t-p} + \epsilon_t\]
Let \(B\) denote the backshift operator. Let \(p(z) = 1 - \sum_{i=1}^p a_i z^i\) denote the AR polynomial. We can write the recursion as
\[p(B)X_t = \epsilon_t\]
A formal solution is given by
\[X_t = p(B)^{-1} \epsilon_t\]
To resolve this formal solution into a well-defined stationary causal solution, we require that \(|z_j| < 1\) for every root \(z_j\) of the AR polynomial \(p(z)\).

Suppose we have data \(X_1, \dots, X_n\). How can we estimate these parameters \(a_1, \dots, a_p\)? The least squares approach is:
\[\min_{a_1, \dots, a_p} \sum_{t=p+1}^n (X_t - (a_1 X_{t-1} + \dots + a_p X_{t-p}))^2\]
Classical time series analysis presumse that \(p\) is fixed but small, so that OLS is a feasible estimation strategy. We can write the normal equations as
\[\begin{bmatrix}
\wh a_1 \\
\vdots \\
\wh a_p
\end{bmatrix}
= (Z^\intercal Z)^{-1} Z^\intercal Y,\]

\[Y = 
\begin{bmatrix}
X_n \\
\vdots \\
X_{p+1}
\end{bmatrix}\]

\[Z =
\begin{bmatrix}
X_{n-1} & \cdots & X_{n-p} \\
\vdots & \ddots & \vdots \\
X_p & \cdots & X_1
\end{bmatrix}\]

We need \(n - p > p\) in order to compute the least squares estimate. Under this regime:
\[\sqrt{n}(\wh a - a) \implies   \N(0, \sigma^2_\epsilon \Gamma_p^{-1})\]
where
\[\Gamma_p =
\begin{bmatrix}
\gamma(0) & \gamma(1) & \cdots & \gamma(p - 1) \\
\gamma(1) & \gamma(0) & \cdots & \gamma(p - 2) \\
\vdots & \vdots & \ddots & \vdots \\
\gamma(p - 1) & \gamma(p - 2) &  \cdots & \gamma(0)
\end{bmatrix}\]
In the high-dimensional regime where \(p \sim cn\), both consistency \& the CLT fail.

If \(p = o(n^{1/3})\), then \(\wh a\) is consistent for \(a\) \citep{Ber74}. Since the dimensionality is growing, we need to be careful to assert the relevant norm in which to measure consistency. The precise result is
\[\max_{1 \leq j \leq p} |\wh a_j - a_j| \cp 0\]
How about other norms?

\subsection*{Vector Autoregression}
Let \(X_t \in \mathbb{R}^p\) for \(t = 1, \dots, p\). The VAR(1) recursion is
\[X_t = AX_{t-1} + \epsilon_t\]
where \(A\) is a \(p \times p\) matrix. For parismony, we might impose structural assumptions on \(A\). We can estimate \(A\) via
\[\min_A \sum_{t=2}^n \|X_t - AX_{t-1}\|_2^2\]
splitting the objective into coordinates \(j = 1, \dots, p\), this becomes
\[\min_{A_{j\cdot}} \sum_{t=2}^n (X_{t,j} - A_{\cdot j} X_{t-1})^2\]
The solution to each of these \(p\) least squares problems \(A_{j\cdot}\) is given by OLS. The design matrix in common across these OLS regressions is
\[\begin{bmatrix}
X_{1,1} & \cdots & X_{1,p} \\
\vdots & \ddots & \vdots \\
X_{n-1, 1} & \cdots & X_{n-1, p}
\end{bmatrix}\]
The row \(A_{j\cdot}\) is obtained by the regression of the following response vector:
\[\begin{bmatrix}
X_{2,j} \\
\vdots \\
X_{n,j}
\end{bmatrix}\]

\subsection*{Dantzig Method}
Right-multiplying the AR(1) recursion by \(X_{t-1}^\intercal\), we have
\[X_t X_{t-1}^\intercal = A X_{t-1} X_{t-1}^\intercal + \epsilon_{t-1} X_{t-1}^\intercal\]
Taking expectations,
\[\E[X_t X_{t-1}^\intercal] = A \E[X_{t-1}X_{t-1}^\intercal] + 0\]
Let
\[\Gamma_0 =
\begin{bmatrix}
E[X_{t,1}^2] & \E[X_{t,1}X_{t,2}] & \cdots & \E[X_{t,1}X_{t,p}] \\
E[X_{t,2}X_{t,1}] & \E[X_{t,2}^2] & \cdots & \E[X_{t,2}X_{t,p}] \\
\vdots & \vdots & \ddots & \vdots \\
E[X_{t,p}X_{t,1}] & \E[X_{t,p}X_{t,2} & \cdots & \E[X_{t,p}^2]
\end{bmatrix}
=
E[X_t X_t^\intercal]\]

\[\Gamma_1 =
\begin{bmatrix}
E[X_{t,1}X_{t-1,1}] & \E[X_{t,1}X_{t-1,2}] & \cdots & \E[X_{t,1}X_{t-1,p}] \\
E[X_{t,2}X_{t-1,1}] & \E[X_{t,2}X_{t-1,2}] & \cdots & \E[X_{t,2}X_{t-1,p}] \\
\vdots & \vdots & \cdots & \vdots \\
E[X_{t,p}X_{t-1,1}] & \E[X_{t,p}X_{t-1,2}] & \cdots & \E[X_{t,p}X_{t-1,p}]
\end{bmatrix}
=
E[X_t X_{t-1}^\intercal]\]

Then we have the population equation
\[\Gamma_1 = A \Gamma_0\]
Let \(\wh\Gamma_1\) and \(\wh\Gamma_0\) be the plug-in estimates. The Dantzig estimate of \(A\) is
\[\min_A \|A\|_1, \quad \text{s.t. } \|\wh\Gamma_1 - A\wh\Gamma_0\|_\infty \leq \lambda\]
The norms \(\|\cdot\|_1\) and \(\|\cdot\|_\infty\) are the vectorized \(\ell_1\) and \(\ell_\infty\) norms. This is a linear program with \(p^2\) inputs \(A_{ij}\). However, this optimization problem decouples into \(p\) subproblems, each with \(p\) inputs, in the following way:

TODO

\subsection*{Inference of Means}
Let \(X_1, \dots, X_n\) be realizations of a stationary process. How can we estimate the marginal mean \(\mu\)? The BLUE estimate can be written as
\[\wt \mu = \sum_{i=1}^n w_i X_i = w^\intercal X\]
Unbiasedness requires that \(\sum_{i=1}^n w_i = 1\). The variance of \(\wt \mu\) is \(w^\intercal \Gamma_n w\). The Lagrangian for this linearly constrained quadratic program is thus
\[w^\intercal \Gamma_n w - \lambda (w^\intercal \mathbf{1} - 1)\]

The stationarity conditions are
\[2\Gamma_n w - \lambda \mathbf{1} = 0 \iff w = \frac{\lambda}{2}\Gamma_n^{-1}\mathbf{1}\]
With the normalization constraint, the solution becomes
\[w = \frac{\Gamma_n^{-1} \mathbf{1}}{\mathbf{1}^\intercal \Gamma_n^{-1} \mathbf{1}}\]
The BLUE is
\[\wt \mu = \frac{\mathbf{1}^\intercal \Gamma_n^{-1} X}{\mathbf{1}^\intercal \Gamma_n^{-1} \mathbf{1}}\]
